\begin{abstract}
배틀로얄 생존은 경기 종료 후의 단일 순위가 아니라 안전 구역 수축에
따라 변하는 위험 과정이다. 본 연구는 공식 PUBG 텔레메트리를 이동하는
팀, 동적 안전/블루존, 공간 이벤트, 팀 관계를 포함하는 시공간
데이터베이스(STDB)로 재구성하고, 재구성한 존 변화 구간에 인덱스를
부여해 팀 단위 다음 페이즈 탈락 위험 예측 문제를 정의한다. 그래프 구성이 가능한
21,824개 경기에서 66차원 팀-페이즈 특징과 근접, 교전, 공동 안전존
압력을 나타내는 3채널 동적 그래프를 구축하였다. 특징만 사용하는
gradient boosting 모델도 AUC 0.765로 강한 성능을 보였지만, 동적
그래프를 사용하는 DynamicSTGNN은 134,245개 정렬 테스트 셀에서
AUC 0.778, AP 0.681까지 향상되었다. 표현 ablation은 텔레메트리
기반 표현이 전략 점수 요약의 AUC 0.649를 0.765까지 끌어올린다는 점을
보였고, 그래프 ablation과 shuffle control은 동적 팀
관계가 특징만으로는 설명되지 않는 작지만 유의한 신호를 제공함을
보였다. 결과적으로 가장 큰 이득은 페이즈 정렬 텔레메트리 표현에서
나오며, 동적 그래프는 그 위에 관계 신호를 추가한다.
\end{abstract}
